\documentclass[11pt]{amsart}
\usepackage{amssymb, latexsym, graphicx, mathrsfs, enumerate, amsmath, amsthm, textcomp, url, tikz-cd, nccmath}
\usepackage[T1]{fontenc}
%\usepackage{mathdots}
\usepackage{verbatim}
\usepackage[all]{xy}
\usepackage{mathtools}

\usepackage{listings}

\lstset{basicstyle=\footnotesize\ttfamily,
        breaklines=true,
        columns=fullflexible}

\usepackage{dsfont}
\renewcommand{\mathbb}{\mathds}

\usepackage{longtable}

%\usepackage[pagebackref=true]{hyperref}
%\renewcommand\backrefxxx[3]{ \hyperlink{page.#1}{$\uparrow$#1} }
%\usepackage{refcheck}

\usepackage[dvipsnames]{xcolor}



% Cyrillic
\DeclareFontEncoding{OT2}{}{} % to enable usage of cyrillic fonts
\DeclareTextFontCommand{\textcyr}{\fontencoding{OT2}
    \fontfamily{wncyr}\fontseries{m}\fontshape{n}\selectfont}
\newcommand{\Sha}{{\!\textcyr{Sh}}}

\theoremstyle{plain}
\newtheorem{theorem}{Theorem}[section]
\newtheorem{conjecture}[theorem]{Conjecture}
\newtheorem{proposition}[theorem]{Proposition}
\newtheorem{lemma}[theorem]{Lemma}
\newtheorem{corollary}[theorem]{Corollary}
\newtheorem{question}[theorem]{Question}
\newtheorem{property}[theorem]{Property}

\theoremstyle{definition}

\newtheorem{example}[theorem]{Example}
\newtheorem{examples}[theorem]{Examples}
\newtheorem{remark}[theorem]{Remark}
\newtheorem{definition}[theorem]{Definition}
\newtheorem{construction}[theorem]{Construction}
\newtheorem{notation}[theorem]{Notation}
%\newtheorem{subsec}[theorem]{}

\newtheorem*{notation*}{Notation}

%\numberwithin{equation}{section}

\usepackage[top=2cm, bottom=3cm, left=3cm, right=3cm]{geometry}
\parskip=4pt


\newcommand{\R}{{\mathbb R}}
\newcommand{\C}{{\mathbb C}}
\newcommand{\A}{{\mathbb A}}
\newcommand{\G}{{\mathbb G}}
\newcommand{\Q}{{\mathbb Q}}
\newcommand{\F}{{\mathbb F}}
\newcommand{\Z}{{\mathbb Z}}

\begin{comment}
\newcommand{\GG}{{\bf G}}
\newcommand{\HH}{{\bf H}}
\newcommand{\NN}{{\bf N}}
\newcommand{\WW}{{\bf W}}
\renewcommand{\AA}{{\bf A}}
\newcommand{\BB}{{\bf B}}
\newcommand{\CC}{{\bf C}}
\newcommand{\LL}{{\bf L}}
\newcommand{\RR}{{\bf L}}
\renewcommand{\SS}{{\bf S}}
\newcommand{\TT}{{\bf T}}
\newcommand{\XX}{{\bf X}}
\newcommand{\YY}{{\bf Y}}
\newcommand{\ZZ}{{\bf Z}}
\end{comment}


\newcommand{\real}{\bf}

\renewcommand{\AA}{{\real A}}
\newcommand{\BB}{{\real B}}
\newcommand{\CC}{{\real C}}
\newcommand{\DD}{{\real D}}
\newcommand{\EE}{{\real E}}
\newcommand{\FF}{{\real F}}
\newcommand{\GG}{{\real G}}
\newcommand{\HH}{{\real H}}
\newcommand{\NN}{{\real N}}
\newcommand{\WW}{{\real W}}
\newcommand{\LL}{{\real L}}
\newcommand{\RR}{{\real R}}
\renewcommand{\SS}{{\real S}}
\newcommand{\TT}{{\real T}}
\newcommand{\UU}{{\real U}}
\newcommand{\XX}{{\real X}}
\newcommand{\YY}{{\real Y}}
\newcommand{\ZZ}{{\real Z}}

\renewcommand{\RR}{{\mathcal R}}
\renewcommand{\SS}{{\mathcal S}}

\newcommand{\X}{{\sf X}}

\newcommand{\gG}{{\scriptscriptstyle \real G}}
\newcommand{\yY}{{\scriptscriptstyle \real Y}}
\newcommand{\xX}{{\scriptscriptstyle \real X}}

\newcommand{\Cc}{{\scriptscriptstyle C}}
\newcommand{\Hh}{{\scriptscriptstyle H}}

\newcommand{\cC}{{\scriptscriptstyle \mathbb C}}
\newcommand{\rR}{{\scriptscriptstyle \mathbb R}}

\newcommand{\cZ}{{\mathcal{Z}}}
\newcommand{\cN}{{\mathcal{N}}}

\newcommand{\bs}{\boldsymbol}

\newcommand{\ag}{{\mathfrak{a}}}
\newcommand{\bg}{{\mathfrak{b}}}
\newcommand{\hg}{{\mathfrak{h}}}
\newcommand{\hR}{\bs{\hg}}
\newcommand{\ig}{{\mathfrak{i}}}
\newcommand{\kg}{\bs{\mathfrak{k}}}
\renewcommand{\ng}{{\mathfrak{n}}}
\newcommand{\tg}{{\mathfrak{t}}}
\newcommand{\tgR}{\bs{\tg}}
\newcommand{\pg}{\bs{\mathfrak{p}}}
\newcommand{\gl}{{\mathfrak{gl}}}
\newcommand{\rg}{{\mathfrak {r}}}
\newcommand{\sg}{{\mathfrak{s}}}
\newcommand{\sgR}{\bs{\sg}}
\newcommand{\zg}{{\mathfrak{z}}}
\newcommand{\zgR}{\bs{\zg}}
\newcommand{\cg}{{\mathfrak{c}}}
\newcommand{\cgR}{\bs{\cg}}
\newcommand{\g}{{\mathfrak{g}}}
\newcommand{\gR}{\bs{\g}}

\newcommand{\Zhat}{\hat{\ZZ}}
%\newcommand{\AAf}{\AA_{\rm f}}
\newcommand{\Fbar}{{\overline{\FF}}}
\newcommand{\Spec}{{\rm Spec}}

\newcommand{\into}{\hookrightarrow}
\newcommand{\rH}{{\rm H}}
\newcommand{\calL}{{\cal L}}
\newcommand{\calV}{{\cal V}}
\newcommand{\calO}{{\cal O}}
\newcommand{\Supp}{{\rm Supp}}

%\newcommand{\bs}{\backslash}
\newcommand{\ti}{\times}
\newcommand{\ol}{\overline}
\newcommand{\wh}{\widehat}
\newcommand{\wt}{\widetilde}
\newcommand{\lto}{\longrightarrow}
\newcommand{\lfrom}{\longleftarrow}
\newcommand{\sm}{{\rm sm}}

\DeclareMathOperator{\Inn}{Inn}
\DeclareMathOperator{\Aut}{Aut}
\DeclareMathOperator{\SAut}{SAut}
\DeclareMathOperator{\Out}{Out}
\DeclareMathOperator{\SOut}{SOut}
\DeclareMathOperator{\inn}{inn}
\DeclareMathOperator{\Lie}{Lie}

\newcommand{\ad}{{\rm ad\hs}}


\newcommand{\End}{{\rm End}}
%\newcommand{\Aut}{{\rm Aut}}
\newcommand{\Hom}{{\rm Hom}}
\newcommand{\Res}{{\rm Res}}
\newcommand{\Sh}{{\rm Sh}}
\newcommand{\Gal}{{\rm Gal}}
%\newcommand{\ad}{{\rm ad}}
\newcommand{\ab}{{\rm ab}}
\newcommand{\der}{{\rm der}}
\newcommand{\nor}{{\rm nor}}
\newcommand{\GL}{{\rm GL}}
\newcommand{\SO}{{\rm SO}}
\newcommand{\bGL}{{\bf GL}}
\newcommand{\MT}{{\rm MT}}
\newcommand{\AM}{{\rm AM}}
\newcommand{\Gm}[2]{{\GG_{{\rm m}{{#1}}}^{{#2}}}}
\newcommand{\ld}{{\langle}}
\newcommand{\rd}{{\rangle}}
\newcommand{\Nm}{{\rm Nm}}
\newcommand{\Cl}{{\rm Cl}}
%\newcommand{\Lie}{{\rm Lie}}
\newcommand{\Ad}{{\rm Ad}}
\newcommand{\diag}{{\rm diag}}

\newcommand{\ov}{\overline}

\newcommand{\isom}{{\stackrel{\sim}{\longrightarrow}}}

\newcommand{\sss}{{\rm ss}}
\newcommand{\ssc}{{\rm sc}}

\newcommand{\isoto}{\overset{\sim}{\longrightarrow}}
\newcommand{\onto}{\twoheadrightarrow}
\newcommand{\labelto}[1]{\xrightarrow{\makebox[1.5em]{\scriptsize ${#1}$}}}
\newcommand{\longisoto}{{\labelto\sim}}

%Small spaces
\newcommand{\hs}{\kern 0.8pt}
\newcommand{\hssh}{\kern 1.2pt}
\newcommand{\hshs}{\kern 1.6pt}
\newcommand{\hssss}{\kern 2.0pt}

\newcommand{\hm}{\kern -0.8pt}
\newcommand{\hmm}{\kern -1.2pt}

\newcommand{\emm}{\bfseries}


%\renewcommand{\H}{{\mathbb H}}
\newcommand{\abar}{{\bar a}}
%\renewcommand{\hbar}{{\bar h}}
\newcommand{\sbar}{{\bar s}}

\newcommand{\upgam}{\hs^\gamma\hm}
\newcommand{\gam}{{\gamma}}
\newcommand{\Gam}{{\Gamma}}


\newcommand{\id}{{\rm id}}
\newcommand{\ups}{{\hs^s\hm}}
\newcommand{\eE}{{\scriptscriptstyle E}}
\newcommand{\spl}{{\rm spl}}


\newcommand{\ff}{{\rm f}}
\newcommand{\uu}{{\rm u}}
\newcommand{\red}{{\rm red}}
\newcommand{\mmu}{{\rm mod-u}}



\newcommand{\im}{{\rm im\,}}
%\newcommand{\cC}{{\scriptscriptstyle \CC}}

\newcommand{\hh}{{\mathfrak h}}

\newcommand{\Stab}{{\rm Stab}}

\newcommand{\SL}{{\rm SL}}
\newcommand{\Oo}{{\mathcal O}}

\newcommand{\GmR}{{\G_{{\rm m},\R}}}
\newcommand{\GmC}{{\G_{{\rm m},\C}}}


\newcommand{\cc}{\raise 1.7pt \hbox{\Tiny{$\bullet$}}}

\newcommand{\Hr}{{\rm H}}
\newcommand{\Zr}{{\rm Z}}
\newcommand{\Br}{{\rm B}}
\newcommand{\Nr}{{\rm N}}

\newcommand{\pn}{{\prime\,0}}
\newcommand{\pf}{{\prime\,\ff}}

\newcommand{\Tm}{{\mathcal T}}


\newcommand{\Ga}{{\Gamma}}

\newcommand{\GGbar}{{\ov\GG}}
\newcommand{\xibar}{{\bar \xi}}
\newcommand{\gbar}{{\bar g}}
\newcommand{\Gbar}{{\ov G}}

\newcommand{\num}{{\scriptscriptstyle\#}}
\newcommand{\sigmai}{{\sigma_i}}
\newcommand{\gi}{{g_i}}
\newcommand{\Xg}{{\mathfrak X}}

%\newcommand{\tT}{{\rm T}}
\newcommand{\ii}{{\boldsymbol{i}}}

\newcommand{\coker}{{\rm coker}}

\newcommand{\SmallMatrix}[1]{{\tiny\arraycolsep=0.4\arraycolsep\ensuremath
{\begin{pmatrix}#1\end{pmatrix}}}}

\newcommand{\vphi}{{\varphi}}

\newcommand{\tl}{\mathfrak{t}}
\newcommand{\ttl}{\boldsymbol{\tl}}

\newcommand{\scong}{{\scriptstyle\cong}}
\newcommand{\mm}{{m}}


\newcommand{\Vssc}{{\vee\hs\ssc}}
\newcommand{\Vad}{{\vee\hs\ad}}

\newcommand{\alg}{{\rm alg}}
%\newcommand{\tT}{{\scriptstyle\rm T}}

\newcommand{\parpr}{{\partial\hs'}}

\newcommand{\ntil}{{\tilde n}}
\newcommand{\obs}{{\rm obs}}

\newcommand{\ug}{{\mathfrak u}}
\newcommand{\llg}{{\mathfrak l}}
\newcommand{\h}{{\mathfrak{h}}}

\newcommand{\bkappa}{{\beta}}
\newcommand{\rr}{{(r)}}

\newcommand{\anti}{{\text{a-r}}}

\newcommand{\V}{{\mathcal V}}

\newcommand{\pia}{{\pi_1^{\rm alg}}}
\newcommand{\Tors}{{\rm Tors}}

\newcommand{\KernelMatrix}{{\tt KernelMatrix}}

\newcommand{\LambdaB}{{\boldsymbol{\Lambda}}}

\newcommand{\Thv}{{\Theta_v}}
\newcommand{\Thvt}{{\Theta_v,\Tors}}
\newcommand{\Tht}{{\Theta,\Tors}}
\newcommand{\ThT}{\Tht}
\newcommand{\vt}{{\vartheta}}

\newcommand{\loc}{{\rm loc}}

\newcommand{\la}{\langle}
\newcommand{\ra}{\rangle}
\newcommand{\smm}{{\sf mm}}

\newcommand{\abLambda}{{{\rm ab}\Lambda}}
\newcommand{\mC}{{\mathcal C}}
\newcommand{\lambdatil}{{\tilde\lambda}}
\newcommand{\cV}{{\mathcal V}}

\newcommand{\vn}{{\varnothing}}
\newcommand{\re}{{\rm re}}

\begin{document}

\title[Relative Tate-Shafarevich kernel]%
{Relative Tate-Shafarevich kernel
of a reductive group over  a global field}

\author{Mikhail Borovoi}


\address{School of Mathematical Sciences,
       Tel Aviv University, 6997801 Tel Aviv, Israel}
\email{borovoi@tauex.tau.ac.il}




\thanks{ The author was partially supported
by the Israel Science Foundation (Grant 1030/22).
}

\keywords{Reductive group,
Galois cohomology, global field, local field, relative Tate-Shafarevich kernel, computer algebra}

\subjclass{%
 20G10% Cohomology theory for linear algebraic groups
, 11E72%   Galois cohomology of linear algebraic groups
, 20G20% Linear algebraic groups over the reals, the complexes, the quaternions
, 20G25% Linear algebraic groups over local fields and their integers
, 20G30% Linear algebraic groups over global fields and their integers
, 68W30% Symbolic computation and algebraic computation
}








\date{\today}

\begin{abstract}

We describe an algorithm and a computer program computing the
relative Tate-Shafarevich kernel $\Sha^S(K,G)$
of a reductive group $G$ over a global field $K$
with respect to a finite set $S$ of places of $K$ containing no archimedean places.
We assume  that we know the algebraic fundamental group $\Lambda=\pi_1^\alg(G)$,
the Galois group $\Theta=\Gal(L/K)$ acting on $\Lambda$, where $L$
is a finite Galois extension of $K$ having no real places,
the decomposition subgroups $\Theta_v\subseteq \Theta$ for all places $v$ in $S$,
and all non-cyclic decomposition subgroups $\Theta_v$ for $v$ outside $S$.
We also compute $\Sha^S(K,G)$  in some cases where $S$ contains archimedean places.



\end{abstract}


\begin{comment}  Abstract for arXiv:

We describe an algorithm and a computer program computing the
relative Tate-Shafarevich kernel $Sha^S(K,G)$
of a reductive group $G$ over a global field $K$
with respect to a finite set $S$ of places of $K$ containing no archimedean places.
We assume  that we know the algebraic fundamental group $\Lambda=\pi_1^{alg}(G)$,
the Galois group $\Theta=\Gal(L/K)$ acting on $\Lambda$, where $L$
is a finite Galois extension of $K$ having no real places,
the decomposition subgroups $\Theta_v\subseteq \Theta$ for all places $v$ in $S$,
and all non-cyclic decomposition subgroups $\Theta_v$ for $v$ outside $S$.
We also compute $Sha^S(K,G)$  in some cases where $S$ contains archimedean places.


\end{comment}



\maketitle

%\tableofcontents


\section{Introduction}

Let $K$ be a global field, that is, a number field or a global function field, and
let $K^s$ be a fixed separable closure of $K$.
We write  $\V(K)$ or $\V_K$ for the set of places of $K$.
Denote by $\V_\infty(K)$ and $\V_f(K)$ the sets of infinite (archimedean)
and finite (non-archimedean) places of $K$, respectively; then $\V(K)=\V_\infty(K)\cup\V_f(K)$.
If $K$ is a function field, then $\V_\infty(K)=\varnothing$ and $\V(K)=\V_f(K)$.

Let $G$ be a (connected) reductive group over $K$.
Write $\Lambda=\pi_1^\alg(G)$ for the algebraic fundamental group of $G$
of \cite[Section 1]{Borovoi-Memoir}.
Then $\Lambda$ is a finitely generated abelian group, on which
the absolute Galois group $\Gal(K^s/K)$ acts (with finite image in $\Aut  \Lambda$).
Let $L$ be a finite Galois extension of $K$ in $K^s$
such that $\Gal(K^s/L)$ acts trivially on $\Lambda$,
and in the case when $K$ is a number field, we choose $L$ to be {\em totally imaginary}.
Then the finite group $\Theta\coloneqq \Gal(L/K)$ naturally acts on $\Lambda$.
(For example, in the case of a function field $K$ or a totally imaginary number field,
we can take for $L$ the finite Galois extension of $K$
corresponding to the kernel of the canonical homomorphism $\Gal(K^s/K)\to\Aut  \Lambda$;
then $\Theta$ acts on $\Lambda$ effectively, that is, with trivial kernel.)
We shall work with the $\Theta$-module $\Lambda$.

For any place $w\in\V(L)$, let $\Theta_w\subseteq\Theta$
denote the corresponding decomposition group,
that is, the stabilizer of $w$ in $\Theta$.
Let $v=v(w)\in\V(K)$ denote the image of $w$ in $\V(K)$.
By abuse of notation, we shall use the notation $\Theta_v$ for $\Theta_w$.
Then $\Theta_v$ is determined by $v$ up to conjugacy in $\Theta$.

Let $S\subset \V(K)$ be a finite set of places.
Write
\[S_\infty=S\cap\V_\infty(K),\quad\  S_f=S\cap \V_f(K),\quad\
       S^\complement=\V(K)\smallsetminus S.\]

Following Sansuc \cite[Notation and conventions]{Sansuc},
we define the relative Tate-Shafarevich kernel with respect to $S$:
\[\Sha^S(K,G)\coloneqq\ker \Big[ H^1(K,G)\to\prod_{v\in S^\complement} H^1(K_v,G)\Big], \]
where $H^1(K,G)$ and $H^1(K_v,G)$ denote the corresponding Galois cohomology sets.
By Borel's theorem (see \cite[page 25]{Borel}) 
this set $\Sha^S(K,G)$ is finite.
In this paper, we describe an algorithm and a Magma program
computing a related finite abelian group $\Sha^S(\Lambda)$.
In the case where $S$ contains no archimedean places (for example, when $K$ is a function field),
there is a canonical bijection $\Sha^1(K,G)\isoto\Sha^S(\Lambda)$.

Note that there are computations of  $\Sha(K,G)\coloneqq\Sha^\varnothing(K,G)$
in \cite[Examples 5.6, 5.7, 5.8]{Sansuc},
\cite[\,(4.2.2)\hs]{Kottwitz-84}, \cite[Corollary 5.13]{Borovoi-Memoir},
and \cite[Theorem 9.4]{CT}
Moreover, Sansuc \cite[p.~15]{Sansuc} computes $\Sha^S(K,B)$ with non-empty $S$
for some {\em finite} abelian $K$-groups $B$;
see \cite[\S 2]{Sansuc}, where he also computes $\Sha_\omega(K,B)=\varinjlim_S \Sha^S(K,B)$.
However, no computations in the literature
of $\Sha^S(K,G)$ for connected reductive $K$-groups $G$ and non-empty $S$ are known to the author.

The plan of the rest of the paper is as follows.
In Section \ref{s:reduction} we reduce computing $\Sha^S(K,G)$
to computing a certain abelian group $\Sha^S(\Lambda)$.
In Section \ref{s:abelian} we give a formula for $\Sha^S(\Lambda)$
in terms of certain  finite abelian groups,
which permits one to compute $\Sha^S(\Lambda)$ using a computer.
In Section \ref{s:algorithm} we describe our algorithm and the Magma program.
In Section \ref{s:example} we consider a certain semisimple groups $G$
over the fields $K=\Q\big(\sqrt7\big)$ and $K=\Q(\alpha)$ where $\alpha$
is a root of the polynomial $f(x)=x^3+x^2-10x-8$.
These groups split over $K(\zeta_8)$ where $\zeta_8$ is a primitive 8th root of unity.
We give  an example of computation
of $\Sha^S(K,G)$ for these $K$ and $G$ using our Magma program.
In the arXiv version of this paper, in Appendix \ref{app:Magma}
we provide a listing of our Magma program.

{\sc Acknowledgements\ } The author thanks Willem A. de Graaf for writing the Magma function
$\tt ComplementIntMat$ used in our Magma program. We thank Taeyeoup Kang for helpful discussions.
The paper and program were proofread by AI Claud Opus 5.5.









\section{Reduction to computing $\Sha^S(\Lambda)$}
\label{s:reduction}

In this section we reduce computing $\Sha^S(K,G)$
to computing a certain abelian group $\Sha^S(\Lambda)$.

Recall that for $v\in \V(K)$, we denote by
$\Thv\subset \Theta$  a decomposition subgroup of $\Theta$ at $v$
(defined up to conjugacy).
We write $\Lambda_\Thv$ for the group of coinvariants of $\Thv$ in $\Lambda$, that is,
\[\Lambda_\Thv=\Lambda/\langle \vartheta\cdot x-x\ | \ x\in \Lambda,\, \vartheta\in\Thv\rangle.\]
We write
\[\Lambda_{\Thv,\Tors}\coloneqq (\Lambda_\Thv)_\Tors\]
for the torsion subgroup of $\Lambda_\Thv$.

Let $\Sigma\subseteq\V_K$ be a subset, finite or infinite.
We write $\Sigma_L$ for the preimage of $\Sigma$ in $\V_L$.
We consider the $\Theta$-modules
\[\Lambda[\Sigma_L]=\Big\{\sum_{w\in\Sigma_L} \lambda_w\cdot w\  \big|\ \lambda_w\in \Lambda\Big\} \]
(the group of formal finite linear combinations), and
\[ \Lambda[\Sigma_L]_0=\Big\{\sum \lambda_w\cdot w\,\in\, \Lambda[\Sigma_L]\ \big|\ \sum \lambda_w=0\Big\}.\]
The Galois group $\Theta$ acts on $\Lambda[\Sigma_L]$ and $\Lambda[\Sigma_L]_0$
diagonally:
\[\vt\cdot \sum \lambda_w\cdot w=
\sum(\vt\cdot \lambda_w)\cdot (\vt\cdot w).
\]
We consider the abelian group
\[ \big(\Lambda[\Sigma_L]_0\big)_{\Theta,\Tors}\hs,\]
the torsion subgroup of the group of coinvariants $\big(\Lambda[\Sigma_L]_0\big)_{\Theta}.$

For a place $v\in\Sigma$, let $\V_L(v)$ denote the preimage of $v$ in $\V_L$ (hence in $\Sigma_L$).
We have a canonical homomorphism
\[\Lambda[\Sigma_L]\to \Lambda[\V_L(v)], \quad\
   \sum_{w\in \Sigma_L}\lambda_w\cdot w\,\mapsto\, \sum_{w\in \V_L(v)} \lambda_w\cdot w.\]
Choose $w\in\V_L(v)$.
Since the $\Theta$-module $\Lambda[\V_L(v)]$ is induced by the $\Theta_w$-module $\Lambda$, by Shapiro's lemma
(see \cite[Proposition 9.76(ii)]{Rotman})
the embedding
\[ \Lambda\into \Lambda[\V_L(v)],\quad\ \lambda\mapsto \lambda\cdot w \]
induces an isomorphism
\[\Lambda_{\Theta_v}\eqqcolon H_0(\Theta_v,\Lambda)\isoto H_0(\Theta,\Lambda[\V_L(v)])\coloneqq \Lambda[\V_L(v)]_\Theta\hs,\]
whence we obtain an isomorphism
\[ \Lambda_\Thvt\isoto \Lambda[\V_L(v)]_\Tht\hs.\]
Thus we obtain a canonical homomorphism
\begin{equation}\label{e:lv}
l_v\colon\,\big(\Lambda[\Sigma_L]_0\big)_{\Theta,\Tors}
    \to \Lambda[\Sigma_L]_{\Theta,\Tors}\to
   \Lambda[\V_L(v)]_{\Theta,\Tors}\overset\sim\longrightarrow \Lambda_{\Thv,\Tors}\hs.
\end{equation}

We refer to \cite[Chapter 5]{BK}
for the definition of the abelian Galois cohomology groups
$H^1_\ab(K,G)$ and $H^1_\ab(K_v,G)$
and the construction of the abelianization maps
\[ \ab\colon H^1(K,G)\to H^1_\ab(K,G)\quad\ \text{and}
    \quad\ \ab_v\colon H^1(K_v,G)\to H^1_\ab(K_v,G).\]

For a {\em real} place $v$ of $K$, since $L$ is totally imaginary, we have
$L_w\simeq\C$ for every place $w$ of $L$ over $v$, hence $\#\Theta_v=2$.
We have the Tate-Nakayama isomorphism
\begin{equation}\label{e:ab-real}
H^1_\ab(K_v,G)\cong H^{-1}(\Theta_v,\Lambda)
\end{equation}
where $H^{-1}$ denotes the corresponding Tate cohomology group;
see \cite[Proposition 8.21]{BT}.
Thus we obtain an embedding
\[H^1_\ab(K_v,G)\cong H^{-1}(\Theta_v,\Lambda)\into \Lambda_{\Theta_v,\Tors}\hs,\]
because by definition
\[H^{-1}(\Theta_v,\Lambda)=\{x\in \Lambda \ |\ x+\vt\cdot x=0\}\hs/\hs\{y-\vt\cdot y\ |\ y\in \Lambda\}\hs\subseteq\hs \Lambda_{\Theta_v,\Tors}\]
where $\Theta_v=\{1,\vt\}$.
For a {\em complex} place $v$ we have the trivial embedding
\[H^1(K_v,G)=H^1_\ab(K_v,G)=\{0\}\into \Lambda_\Thvt=\Lambda_\Tors\hs.\]


For a {\em finite} place $v$ of $K$, we have a canonical isomorphism
$$H^1_\ab(K_v,G)\cong \Lambda_{\Theta_v,\Tors};$$
see \cite[Theorem 5.5.2(1)]{BK}.
For our global field $K$, we have a canonical map
\begin{equation}\label{e:ab}
\abLambda\colon H^1(K,G)\to H^1_\ab(K,G)\to \big(\Lambda[\V_L]_0\big)_\Tht\hs;
\end{equation}
see \cite[Theorem 5.5.4(1)]{BK}.

Set
\[\LambdaB_{\Sigma_\infty}=\bigoplus_{v\in\Sigma_\infty} \Lambda_{\Thv,\Tors}\]
where $\Sigma_\infty=\Sigma\cap\V_\infty(K)$.
We have a canonical homomorphism
\[\bigoplus_{v\in\Sigma_\infty} l_v\colon\,\big(\Lambda[\Sigma_L]_0\big)_{\Theta,\Tors}\to
      \bigoplus_{v\in\Sigma_\infty} \Lambda_{\Thv,\Tors}=\LambdaB_{\Sigma_\infty}\hs.\]
On the other hand, for any $v\in\Sigma_\infty$\hs, using \eqref{e:ab-real} we obtain a canonical map
\[ H^1(K_v,G)\to H^1_\ab(K_v,G)\into \Lambda_{\Thv,\Tors}\hs,\]
whence we obtain a canonical map
\[ \ab\Lambda_{\Sigma_\infty}\colon\prod_{v\in \Sigma_\infty} H^1(K_v,G)\to \LambdaB_{\Sigma_\infty}\hs.\]

Let $S$ be a subset of $\V_K$ and write $S_L$ for the preimage of $S$ in $\V_L$\hs.
Write $S^\complement=\V_K\smallsetminus S$ for the complement of $S$ in $\V_K$\hs,
and $S_L^\complement=\V_L\smallsetminus S_L$ for the complement of $S_L$ in $\V_L$.
We define
\begin{align*}
\Sha^S_\ab(K,G)&=\ker \Big[ H^1_\ab(K,G)\to\prod_{v\in S^\complement} H^1_\ab(K_v,G)\Big],\\ \notag
\Sha^S(\Lambda)&=\ker \Big[ \big(\Lambda[\V_L]_0\big)_{\Theta,\Tors}
     \to \Lambda\big[S_L^\complement\big]_{\Theta,\Tors}\Big]\\
   &= \ker \Big[\bigoplus_{v\in S^\complement}l_v\colon\hs  \big(\Lambda[\V_L]_0\big)_{\Theta,\Tors}\to \bigoplus_{v\in S^\complement} \Lambda_\Thvt\,\Big],
\end{align*}
where
\[\Lambda\big[S_L^\complement\big]_{\Theta,\Tors}
   \cong \hs \bigoplus_{v\in S^\complement} \Lambda_\Thvt\quad\  \text{because} \quad\ \Lambda\big[S_L^\complement\big]=\bigoplus_{v\in S^\complement}\Lambda[\V_L(v)].\]
The map \eqref{e:ab} induces an abelianization map
\begin{equation}\label{e:ab-Sha}
\abLambda^S\colon \Sha^S(K,G)\to \Sha^S_\ab(K,G)\to \Sha^S(\Lambda).
\end{equation}
Moreover, we have a natural map
\[ l_{S_\infty}=\prod_{v\in S_\infty}l_v\colon\hs \Sha^S(\Lambda)\to \bigoplus_{v\in S_\infty} \Lambda_\Thvt\eqqcolon\LambdaB_{S_\infty}\hs.\]

\begin{theorem}\label{t:Sha}
Let $S$ be a subset of $\V_K$\hs, finite or infinite.
The following diagram is  Cartesian:
\begin{equation}\label{e:Cart-S}
\begin{aligned}
\xymatrix@C=14mm{
\Sha^S(K,G)\ar[r]^-{\abLambda^S} \ar[d]_-{\loc_{S_\infty}}     &\Sha^S(\Lambda)\ar[d]^-{l_{S_\infty}} \\
\prod\limits_{v\in S_\infty}\!\! H^1(K_v,G)\ar[r]^-{\ab\Lambda_{S_\infty}}  &\LambdaB_{S_\infty}\hskip-5mm
}
\end{aligned}
\end{equation}
\end{theorem}

\begin{proof}
Diagram \eqref{e:Cart-S} naturally embeds into the following diagram:
\begin{equation}\label{e:Cart}
\begin{aligned}
\xymatrix@C=14mm{
H^1(K,G)\ar[r]^-{\abLambda} \ar[d]_-{\loc_\infty}  &\big(\Lambda[\V_L]_0\big)_\Tht\hskip-7mm\ar[d]^-{l_\infty} \\
\prod\limits_{v\in \V_\infty(K)}\!\! H^1(K_v,G)\ar[r]^-{\ab\Lambda_{\V_\infty(K)}}
       &\LambdaB_{\V_\infty(K)}\hskip-8.5mm
}
\end{aligned}
\end{equation}
which is Cartesian by \cite[Theorem 5.5.10]{BK}.
It follows that diagram \eqref{e:Cart-S} is commutative.
We show that it is Cartesian.
We must prove that if elements $x_{S_\infty}\in \prod_{v\in S_\infty} H^1(K_v,G)$ and
$x_\Lambda\in \Sha^S(\Lambda)$ have the same image in
$\LambdaB_{S_\infty}$,
then they are the images of a unique element
$x\in\Sha^S(K,G)$.

Let
\[ x_\infty\in  \prod\limits_{v\in \V_\infty}\!\! H^1(K_v,G)\]
be the extension of $x_{S_\infty}$ by the trivial classes
at $v\in\V_\infty(K)\smallsetminus S_\infty$.
Then $x_{\infty}$ and $x_\Lambda$ have the same image in $\LambdaB_{\V_\infty(K)}$.
For $v\in S_\infty$ this is the hypothesis.
For each archimedean $v\notin S$ we have $l_v(x_\Lambda)=0$,
which holds because $x_\Lambda\in\Sha^S(\Lambda)$.
Since diagram \eqref{e:Cart} is Cartesian,
the elements $x_{\infty}$ and $x_\Lambda$ are the images of a unique element $x\in H^1(K,G)$,
and it remains to show that $x\in\Sha^S(K,G)$.
By construction, for each {\em infinite} $v\in \V_\infty(K)\smallsetminus S_\infty$\hs,
the image of $x$ in  $H^1(K_v,G)$ is trivial.
Moreover, for each {\em finite} place $v\in\V_f(K)\smallsetminus S_f$\hs,
by \cite[Proposition 4.6.2 together with Theorems 5.5.2(1) and 5.5.4(1)]{BK}
the diagram below is commutative, and
the bottom horizontal arrow in this commutative diagram is bijective:
\[
\xymatrix@C=14mm{
H^1(K,G)\ar[r]^-\abLambda\ar[d]_-{\loc_v} & \big(\Lambda[\V_L]_0\big)_\Tht\hskip-7mm\ar[d]^-{l_v}\\
H^1(K_v,G) \ar[r]^-{\ab\Lambda_v}  &\Lambda_\Thvt\hskip-10mm
}
\]
and therefore for $l_v$ as in \eqref{e:lv}, from $l_v(\abLambda(x))=0$ we conclude that $\loc_v(x)=1$.
Thus for all
$$v\in(\V_\infty(K)\smallsetminus S_\infty)\cup(\V_f(K)\smallsetminus S_f)=S^\complement$$
we have $\loc_v(x)=1$, that is, $x\in\Sha^S(K,G)$.
We see that diagram \eqref{e:Cart-S} is Cartesian,
which completes the proof of the theorem.
\end{proof}

\begin{corollary}
If $S$ contains no infinite (archimedean)  places, that is, $S_\infty=\varnothing$,
then the abelianization map $\abLambda^S$ of  \eqref{e:ab-Sha} is bijective.
\end{corollary}

This corollary defines a canonical abelian group structure
on the pointed set $\Sha^S(K,G)$ in the case when $S$ contains no infinite places.
It is easy to see that this abelian group structure is functorial in $G$ and $S$.

\begin{corollary}
If $K$ is a global function field, then the abelianization map
$\abLambda^S$ of  \eqref{e:ab-Sha} is bijective.
\end{corollary}

\begin{proof}
Indeed, then $S$ contains no archimedean places.
\end{proof}






\section{$\Sha^S(\Lambda)$ in terms of finite abelian groups}
\label{s:abelian}

By Theorem \ref{t:Sha}, in order to compute $\Sha^S(K,G)$ for a finite subset $S\subset \V(K)$, it suffices to compute the finite abelian group $\Sha^S(\Lambda)$ and, for each {\em real} place $v\in S_\infty$, the finite pointed set
$H^1(K_v,G)$ for $v\in S_\infty\subseteq\V_\infty(K)$,
together with the maps
\[\abLambda_v\colon\,  H^1(K_v,G)\to H^1_\ab(K_v,G)\cong H^{-1}(\Theta_v,\Lambda)\into  \Lambda_\Thvt\]
and
\[ l_v\colon \Sha^S(\Lambda)\to \Lambda_\Thvt\]
The real cohomology sets $H^1(K_v,G)$ can be computed using the computer program of \cite{BdG}.
This program gives explicit cocycles in $Z^1(K_v,G)$,
and so it is not difficult to compute their images in
$\Lambda_\Thvt$ under $\abLambda_v$; we shall describe this in detail elsewhere.

It remains to compute the finite abelian group $\Sha^S(\Lambda)$
together with the homomorphisms $l_v$ for $v\in S_\infty$.
In this section,  we give a formula for $\Sha^S(\Lambda)$
in terms of certain  finite abelian groups,
and  formulas for $l_v$.
This permits one to compute $\Sha^S(\Lambda)$ and the homomorphisms $l_v$ using a computer.

For a finite Galois extension of global fields $L/K$ with Galois group $\Theta=\Gal(L/K)$,
and  a finite or infinite  set of places $S\subset \V_K$,
we consider a finite or infinite set $\Sigma\subseteq\V(K)$ containing $S$.
We denote by $S_L$ and $\Sigma_L$ the preimages of  $S$ and $\Sigma$ in $\V(L)$,
respectively.

\begin{definition}\label{d:sl}
For $S\subseteq\Sigma\subseteq \V(K)$,
we say that  $\Sigma$ is {\em sufficiently large for $S$}
if for any place $w\in \Sigma_L^\complement$
there exists a  place $w_{_\Sigma}\in \Sigma_L\smallsetminus S_L$
such that
\begin{equation*}
\Theta_{w_{_\Sigma}}\supseteq \Theta_w\hs.
\end{equation*}
\end{definition}

\begin{lemma} \label{l:S-Sigma}
For any finite set $S\subset \V_K$,
there exists a {\emm finite} set $\Sigma\subset \V(K)$ containing $S$
that is sufficiently large for $S$.
\end{lemma}

\begin{proof}
Let $\V_{\rm nc}(L/K)\subset\V(K)$ denote the set of places of $K$
with non-cyclic decomposition group in $L$.
At each place $v\in\V_{\rm nc}(L/K)$, the extension $L/K$ ramifies,
and therefore there are only finitely many such places.
Thus we may choose $\Sigma$ containing the finite sets $S$ and  $\V_{\rm nc}(L/K)$.
Since our Galois group $\Theta=\Gal(L/K)$ is finite,
there are finitely many conjugacy classes $\mC$ of maximal cyclic subgroups $C\subseteq\Theta$.
For each such conjugacy class $\mC$ we choose $C\in\mC$.
Since $S_L$ is finite, by the Chebotarev density theorem
we can choose a finite place $w_{C}\in (\V_L\smallsetminus S_L)_f$ with decomposition group $C$.
Let $v_\mC \in (\V_K\smallsetminus S)_f$ denote the image of $w_C$.
We may choose $\Sigma$ containing the finite set of these $v_\mC$.


Now for $\Sigma$ as above and for any $w\in \Sigma_L^\complement$,
the decomposition group $\Theta_w$ of $w$ is cyclic,
and hence $\Theta_w\subseteq C_0$ for some maximal cyclic subgroup $C_0$ of $\Theta$.
By construction, there exists a finite place $w_C\in \Sigma_L\smallsetminus S_L$
with decomposition group $\Theta_{w_C}=\vt C_0\vt^{-1}$ for some $\vt\in\Theta$.
Set $w'= \vt^{-1}\cdot w_C$;
then $\Theta_{w'} =C_0$.
Since both $\Sigma_L$ and $S_L$ are $\Theta$-invariant, we have
$w'\in\Sigma_L\smallsetminus S_L$.
We obtain that $\Theta_w\subseteq C_0=\Theta_{w'}$ for $w'\in \Sigma_L\smallsetminus S_L$\hs.
Thus  $\Sigma$ is sufficiently large for $S$ and finite, as desired.
\end{proof}


\begin{construction}
Let $S\subset \V_K$ be a  set of places of $K$,
and let $\Sigma\subseteq\V_K$ be a sufficiently large set for $S$, finite or infinite.
We construct a $\Theta$-equivariant map
\begin{equation*}
\phi_\Sigma\colon \Sigma_L^\complement \to\Sigma_L\smallsetminus S_L
\end{equation*}
as follows.
For any $v\in \Sigma^\complement$
we choose a place $w\in \Sigma_L^\complement$ over $v$.
Then there exists a place $w_\Sigma\in \Sigma_L\smallsetminus S_L$
with $\Theta_{w_\Sigma}\supseteq \Theta_w$.
We define the map $\phi_\Sigma$ on the orbit $\V_L(v)=\Theta\cdot w$
by $\phi_\Sigma(\vartheta\cdot w)=\vartheta\cdot w_\Sigma$ for $\vartheta\in\Theta$.
Since $\Theta_{w_\Sigma}\supseteq \Theta_w$\hs,
this map is well defined.

We extend the map $\phi_\Sigma$ to a map
\begin{equation}\label{e:phi-ext}
\wh\phi_\Sigma\colon \V_L \to \Sigma_L
\end{equation}
by putting $\wh\phi_\Sigma(w)=w$ for $w\in \Sigma_L$. Let
\[\iota_\Sigma\colon \Sigma_L\into\V_L\]
denote the inclusion map.
Then
\begin{equation*}
\wh\phi_\Sigma\circ\iota=\id.
\end{equation*}
By abuse of notation, we write $\phi_\Sigma$ for $\wh\phi_\Sigma$.
Then the map $\phi_\Sigma\colon \V_L\to\Sigma_L$ is surjective.
\end{construction}

For a finite or infinite set $\Sigma$ satisfying $S\subseteq\Sigma\subseteq\V_K$,
we define
\begin{multline*}
\ker_\Sigma^S(\Lambda)=\ker \Big[\big(\Lambda[\Sigma_L]_0\big)_\Tht\to \prod_{v\in \Sigma\smallsetminus S}\Lambda_\Thvt\Big]\\
   =\ker \Big[\big(\Lambda[\Sigma_L]_0\big)_\Tht\to \Lambda[\Sigma_L\smallsetminus S_L]_\Tht\Big]
\end{multline*}
where the homomorphism $\big(\Lambda[\Sigma_L]_0\big)_\Tht\to \Lambda[\Sigma_L\smallsetminus S_L]_\Tht$ is given by restriction from $\Sigma_L$ to $\Sigma_L\smallsetminus S_L$\hs:
\[ \bigg[\,\sum_{w\in \Sigma_L} \lambda_w\cdot w\hs\bigg]\longmapsto
    \bigg[\,\sum_{w\in \Sigma_L\smallsetminus S_L}\!\!\!\! \lambda_w\cdot w\hs\bigg].\]
Then
\[ \ker^S_{\V(K)}(\Lambda)=\Sha^S(\Lambda).\]


\begin{theorem}\label{t:Sigma'}
Let $L/K$ be a finite Galois extension of global fields with Galois group $\Theta=\Gal(L/K)$.
Let
\[ S\subseteq\Sigma\subseteq\Sigma'\subseteq\V_K\]
be subsets of $\V_K$ where $\Sigma$ is sufficiently large for $S$.
Then the natural homomorphism
\[\ker^S_\Sigma(\Lambda)\to\ker^S_{\Sigma'}(\Lambda)\]
is an isomorphism.
\end{theorem}

\begin{corollary}\label{c:Sigma}
The natural homomorphism
$\ker^S_\Sigma(\Lambda)\to\Sha^S(\Lambda)$
is an isomorphism.
\end{corollary}

\begin{proof}[Proof of the corollary] We take $\Sigma'=\V_K$.\end{proof}

\begin{proof}[Proof of Theorem \ref{t:Sigma'}]
We denote by $\phi'_\Sigma\colon \Sigma'_L\to\Sigma_L$ the restriction to $\Sigma'$
of the map $\phi_\Sigma \colon \V_L\to \Sigma_L$ of \eqref{e:phi-ext}.
To simplify the notations, we write
$\phi$ for $\phi'_\Sigma$.
From our construction of the map $\phi$, it follows that
$\phi(\Sigma'_L\smallsetminus S_L)\subseteq \Sigma_L\smallsetminus S_L$.
From the following commutative diagram:
\[
\xymatrix{
\Lambda[\Sigma'_L]_0\ar[r]\ar[d]^{\phi_*} &\Lambda[\Sigma'_L\smallsetminus S_L]\ar[d]^{\phi_*}\\
\Lambda[\Sigma_L]_0\ar[r]                &\Lambda[\Sigma_L\smallsetminus S_L]
}
\]
we obtain a commutative diagram with exact rows
\[
\xymatrix{
0\ar[r] &\ker^S_{\Sigma'}\ar@{-->}[d]^{\phi_*}\ar[r]
        &\big(\Lambda[\Sigma'_L]_0\big)_\Tht\ar[r]\ar[d]^{\phi_*}
       &\Lambda[\Sigma'_L\smallsetminus S_L]_\Tht\ar[d]^{\phi_*}\\
0\ar[r] &\ker^S_{\Sigma}\ar[r]&\big(\Lambda[\Sigma_L]_0\big)_\Tht \ar[r]
       &\Lambda[\Sigma_L\smallsetminus S_L]_\Tht\,.
}
\]
Let $\iota\colon \Sigma_L\into \Sigma'_L$ denote the inclusion homomorphism.
Since by construction we have
\[\phi\circ\iota=\id\colon \Sigma_L\into\Sigma'_L\onto\Sigma_L,\]
we see that
\[\phi_*\circ \iota_*=\id\colon\hs \ker^S_\Sigma(\Lambda)\to\ker^S_{\Sigma'}(\Lambda)\to \ker^S_\Sigma(\Lambda).\]
Below we show that $\iota_*\circ\phi_*=\id$;
this will prove that $\iota_*$ is an isomorphism.

Let
\[
\lambdatil=\sum_{w\in \Sigma'_L} \lambda_w\cdot w \in \Lambda[\Sigma'_L]_0
\]
with $[\lambdatil]\in\ker^S_{\Sigma'}$.
Then $\sum_{w\in \Sigma_L'}\lambda_w=0$ and for each $v\in \Sigma'\smallsetminus S$ we have
\[\sum_{w\in \V_L(v)} \lambda_w\cdot w=\sum_{\vt\in\Theta} (\vt-1)\cdot \alpha_{v,\vt}\]
for some $\alpha_{v,\vt}\in \Lambda[\V_L(v)]$.
In particular, this holds for all $v\in\Sigma'\smallsetminus \Sigma$.

We have
\begin{align*}
(\iota_*\circ\phi_*)(\lambdatil)-\lambdatil
&=\sum_{v\in\Sigma'}\bigg( \phi\Big(\sum_{w\in\V_L(v)}\!\!\lambda_w\cdot w\Big)
   -\sum_{w\in\V_L(v)}\!\!\lambda_w\cdot w\bigg)\\
&=\sum_{v\in\Sigma'\smallsetminus\Sigma}\bigg( \phi\Big(\sum_{w\in\V_L(v)}\!\!\lambda_w\cdot w\Big)
   -\sum_{w\in\V_L(v)}\!\!\lambda_w\cdot w\bigg)
\end{align*}
where we may take the sum over $\Sigma'\smallsetminus\Sigma$
because for $v\in\Sigma$ and $w\in\V_L(v)$ we have $\phi(w)=w$.
It follows that
\begin{align*}
(\iota_*\circ\phi_*)(\lambdatil)-\lambdatil
&=\sum_{v\in\Sigma'\smallsetminus\Sigma}\phi\Big(\sum_{\vt\in\Theta}(\vt-1)\alpha_{v,\vt}\Big)
   -\sum_{v\in\Sigma'\smallsetminus\Sigma\ }\sum_{\vt\in\Theta}(\vt-1)\alpha_{v,\vt}\\
&=   \sum_{v\in\Sigma'\smallsetminus\Sigma}
    \ \sum_{\vt\in\Theta}(\vt-1)\cdot\big(\phi_*(\alpha_{v,\vt})-\alpha_{v,\vt}\big).
    \end{align*}
Write
\[\alpha_{v,\vt}=\sum_{w\in \V_L(v)}\!\! a_w\cdot w\in \Lambda[\V_L(v)]\subset \Lambda[\Sigma'_L] \]
and set
\[\beta_{v,\vartheta}=\phi_*(\alpha_{v,\vt})-\alpha_{v,\vt}=
         \sum_{w\in \V_L(v)}\!\! a_w\cdot \big(\phi(w)-w\big)\in \Lambda[\Sigma'_L]_0\hs.\]
Note that while $\alpha_{v,\vt}\in \Lambda[\Sigma'_L]$, we have
$\beta_{v,\vartheta}\in \Lambda[\Sigma'_L]_0$\hs.
Then
\[(\iota_*\circ\phi_*)(\lambdatil)-\lambdatil=
     \sum_{v\in\Sigma'\smallsetminus\Sigma}
     \ \sum_{\vt\in\Theta}(\vt-1)\cdot\beta_{v,\vt}
     =\ \sum_{\vt\in\Theta}\bigg((\vt-1)
     \cdot\!\!\!\!\sum_{v\in\Sigma'\smallsetminus\Sigma}\!\!\!\beta_{v,\vt}\bigg)
     \]
 with $\beta_{v,\vt}\in \Lambda[\Sigma'_L]_0$\hs, whence
\[ [(\iota_*\circ\phi_*)(\lambdatil)]-[\lambdatil]=0\in \big(\Lambda[\Sigma'_L]_0\big)_\Tht\]
and $(\iota_*\circ\phi_*)([\lambdatil])=[\lambdatil]$.
Thus
\[\iota_*\circ \phi_*=\id\colon\hs \ker^S_{\Sigma'}(\Lambda)\to\ker^S_{\Sigma}(\Lambda)\to \ker^S_{\Sigma'}(\Lambda)\]
and $\iota_*$ is an isomorphism, which completes the proof of the theorem.
\end{proof}

\begin{proposition}
Let $S_1\subset \V_K$ and $S_2\subset \V_K$ be two finite subsets
satisfying  $S_1\subseteq S_2\subset\V_K$.
Let $\Sigma\subset \V_K$ be a sufficiently large finite subset
for $S_2$ (and hence for  $S_1$).
Then the following diagram with natural arrows commutes:
\[
\xymatrix{
\ker_\Sigma^{S_1}(\Lambda)\ar[r]^-\sim\ar[d]  & \Sha^{S_1}(\Lambda)\ar[d]\\
\ker_\Sigma^{S_2}(\Lambda)\ar[r]^-\sim  & \Sha^{S_2}(\Lambda)
}
\]
\end{proposition}
The proposition follows immediately from the naturality of $i_*$.












\newcommand{\vk}{{\varkappa}}


\section{Algorithm}
\label{s:algorithm}


Let the finite group $\Theta=\Gal(L/K)$ be as above, and denote by
\[\Lambda=\pi_1^\alg(G)\]
the algebraic fundamental group of \cite{Borovoi-Memoir}.
We choose a finite system of generators  $e_1,\dots,e_\smm$ of $\Lambda$.
We have a homomorphism
\begin{equation*}
\vk\colon \Z^{\smm}\to\Lambda.
\end{equation*}
Then $\ker\vk$ is a  finitely generated subgroup of $\Z^{\smm}$;
let $\tt N$ be an $\tt nn\times\smm$ matrix whose rows generate $\ker\vk$.
We denote  $\tt M=I_\smm$ (the identity matrix); then $\Lambda=\tt \la M\ra/\la N\ra$.

The finite group $\Theta$ acts on the finitely generated group $\Lambda$.
To any $\vt\in\Theta$ we have an $\smm\times\smm$ integer matrix (defined not uniquely)
$\rho(\vt)=\big(\rho(\vt)_i^j\big)_{1\le i,j\le\smm}$ such that
\[\vt(e_i)=\sum_j\rho(\vt)_i^j\hs e_j\hs\]
where $\rho(\vt)_i^j$is the entry in row $j$, column $i$.
In the program, we {\em define} $\Lambda$ by  integers $\smm$ and $\tt nn$
and an $\tt nn\times\smm$  matrix $\tt N$ as $\Lambda=\tt\la M\ra/\la N\ra$.
We define an action of $\Theta$ on $\Lambda$ by a family $\big(\rho(\vt)\big)_{\vt\in\Theta}$
of $\smm\times\smm$ integer matrices.
Using our function {\tt IsAction}, we check that these data
indeed define a left action of $\Theta$ on $\Lambda$, that is,
\[(\rho(1)-I_\smm)\cdot\tt \Z^\smm\subseteq\la N\ra\quad\text{and}
\quad \big(\rho(\vt_1\vt_2)-\rho(\vt_1)\rho(\vt_2)\big)\cdot \tt \Z^\smm\subseteq\la N\ra.\]

We assume that we are given a set $S$, and we choose a sufficiently large set $\Sigma=S\cup (\Sigma\smallsetminus S)$ (Definition \ref{d:sl}).
In terms of places of $K$ this means:
\begin{equation}
\begin{aligned}\label{e:v-Sigma}
&\textit{for any decomposition subgroup $\Theta_v\subseteq \Theta$
for $v\in \V(K)\smallsetminus \Sigma$, there exists} \\
&\textit{$v_\Sigma\in \Sigma\smallsetminus S$
with $\Theta_{v_\Sigma}\supseteq \Theta_v$ up to conjugacy.}
\end{aligned}
\end{equation}
We assume that we know the decomposition subgroups $\Theta_v$ for all $v\in \Sigma$.
We construct a $\Theta$-set $\Sigma_L$ by
\[ \Sigma_L=\bigsqcup_{v\in\Sigma} \Theta/\Theta_v\hs.\]
Consider the group $\Lambda[\Sigma_L]$ on which the  finite group $\Theta=\Gal(L/K)$ and the group ring
$\Z[\Theta]$ act.
We wish to compute
\[ \ker_\Sigma^S(\Lambda)\coloneqq
\ker\Big[ \big(\Lambda[\Sigma_L]_0\big)_\ThT\to \Lambda[\Sigma_L\smallsetminus S_L]_\Theta\Big].\]
We also consider the ``extended kernel''
\[ {\rm Eker}_\Sigma^S(\Lambda)\coloneqq
\ker\Big[ (\Lambda[\Sigma_L]_0)_\Theta\to \Lambda[\Sigma_L\smallsetminus S_L]_\Theta\Big];\]
then $\ker_\Sigma^S(\Lambda)=\big({\rm Eker}_\Sigma^S(\Lambda)\big)_{\rm Tors}\hs.$

By abuse of notation, we write $\Sigma$ for $\Sigma_L$ and $S$ for $S_L$.
We have
\[{\rm Eker}_\Sigma^S(\Lambda)=\ker \Bigg[\frac{\Lambda[\Sigma]_0}{I_\Theta\, \Lambda[\Sigma]_0}
\,\to\,
\frac{\Lambda[\Sigma\smallsetminus S]}{ I_\Theta\, \Lambda[\Sigma\smallsetminus S]}\Bigg] \]
where
\[ \Lambda[\Sigma]_0=\Big\{\sum \lambda_w\cdot w\in\Lambda[\Sigma]\ \big|\ \sum\lambda_w=0\Big\},\]
and $I_\Theta\subset \Z[\Theta]$
is the ideal generated by the elements $\vt-1$ for $\vt\in\Theta$.
We see that
\[{\rm Eker}_\Sigma^S(\Lambda)= \frac
{\Lambda[\Sigma]_0\cap \big(\Lambda[S]\oplus I_\Theta\, \Lambda[\Sigma\smallsetminus S]\big)}
{I_\Theta\,\Lambda[\Sigma]_0}.
\]

We pass to the preimages in $\la M\ra[\Sigma]$  of the three groups
\[\Lambda[\Sigma]_0,\quad I_\Theta\,\Lambda[\Sigma]_0,\quad
\Lambda[S]\oplus I_\Theta\,\Lambda[\Sigma\smallsetminus S] \]
under the canonical homomorphism $\la M\ra[\Sigma] \to \Lambda[\Sigma]$.
These preimages are:
\begin{gather*}
\la M\ra[\Sigma]_0+ \la N\ra[\Sigma],
\quad
I_\Theta\,\la M\ra[\Sigma]_0 + \la N\ra[\Sigma],
\quad
\la M\ra[S]\,\oplus\, I_\Theta\,\la M\ra[\Sigma\smallsetminus S] + \la N\ra[\Sigma].
\end{gather*}
We obtain that
\[ {\rm Eker}_\Sigma^S(\Lambda)\,\cong\,\frac
{\big(\,\la M\ra[\Sigma]_0 +\la N\ra[\Sigma]\,\big)\,\cap\,
\big(\,\la M\ra[S]\oplus I_\Theta\,\la M\ra[\Sigma\smallsetminus S]
+\la N\ra[\Sigma]\,\big)}
{ I_\Theta\,\la M\ra[\Sigma]_0+\la N\ra[\Sigma]}\,.
\]

It is easy to compute  three matrices
$\tt Z0, C0$, and $\tt CC$ in Magma such that
\begin{gather*}
\la{\tt Z0}\ra=\la M\ra[\Sigma]_0,\quad \la{\tt C0}\ra=  I_\Theta\,\la M\ra[\Sigma]_0,\quad
\la{\tt CC}\ra= \la M\ra[S]\,\oplus\, I_\Theta\,\la M\ra[\Sigma\smallsetminus S].
\end{gather*}
Let $\tt NN$ denote the {\em diagonal join}
of $\#\Sigma_L$ copies of $N$;
then $\la N\ra[\Sigma]={\tt\la NN\ra}$.
We compute three matrices
\begin{gather*}\tt Z0N:=VerticalJoin(Z0,NN),\quad
C0N:=VerticalJoin(C0,NN),\\
\tt CCN:=VerticalJoin(CC,NN).
\end{gather*}
Then
\begin{gather*}
{\tt \la Z0N\ra}=\la M\ra[\Sigma]_0 + \la N\ra[\Sigma],\quad
\la{\tt C0N}\ra=  I_\Theta\,\la M\ra[\Sigma]_0 + \la N\ra[\Sigma],\\
\la{\tt CCN}\ra= \la M\ra[S]\,\oplus\, I_\Theta\,\la M\ra[\Sigma\smallsetminus S]
+ \la N\ra[\Sigma].
\end{gather*}
Using our Magma function Meet, we compute a matrix
\[\tt mZ0CCN:=Meet(Z0N, CCN)\quad\
\text{such that}\quad\
\tt \la mZ0CCN\ra=\la Z0N\ra\cap\la CCN\ra .
\]
Using our Magma function ${\tt ComplementIntMat}$ (written by Willem A. de Graaf),
we compute
\[\tt comp, bas, x := ComplementIntMat(\ mZ0CCN,\  C0N\ ).\]
Here $\tt bas$ is a matrix with $n_{\tt bas}$ rows $\tt bas[i]$
and $\tt x$ is a sequence of positive integers $\tt x[i]$ of length $n_{\tt bas}$
such that  ${\tt x[i]* bas[i]}, \ i=1,\dots,\ n_{\tt bas}\hs$,
constitute a basis of $\tt \la C0N\ra$
and that  the vectors $\tt bas[i]$ belong to $\tt \la mZ0CCN\ra$
and represent a set of generators of the finite abelian group
\begin{equation*}
\tt \big(\la mZ0CCN\ra/  \la C0N\ra\big)_\Tors\hs ,
\end{equation*}
with each class $\tt [bas[i]]$ of order $\tt x[i]$.
Furthermore, we have an isomorphism of finite abelian groups
\[\tt \bigoplus_i \Z/x[i]\Z
\cong(\,{\tt\la mZ0CCN\ra\, /\,  \la C0N \ra}\, )_{\rm Tors}\]
sending $[1]\in \tt \Z/x[i]\Z$ to the class $\tt [bas[i]\hs]$ of $\tt bas[i]$.
It follows that
\[\Sha^S(\Lambda)\cong \ker_\Sigma^S(\Lambda)
\cong (\,{\tt\la mZ0CCN\ra\, /\,  \la C0N \ra}\, )_{\rm Tors}\cong \bigoplus_i \Z/x[i]\Z,\]
which computes $\Sha^S(\Lambda)$.

It remains to compute the homomorphisms
\begin{equation}\label{e:lvv} l_v\colon \Sha^S(\Lambda)\cong\ker_\Sigma^S(\Lambda)\to \Lambda_\Thvt
\end{equation}
for real places $v$ in $S$.
It suffices to compute the images
$l_v\tt\big[ bas[i]\big]$ of the generators
$\tt \big[ bas[i]\big]$.
We may regard each  $\tt bas[i]$ as an element $\la M\ra[\Sigma]$,
say, ${\tt bas[i]}=\sum_{w\in\Sigma} m_w\cdot w$.

Let $v\in S_\infty\subset \Sigma$. Choose a place  $w_0$ of $L$ over $v$.
Recall that $\V_L(v)$ denotes the preimage of $v$ in $\V_L$.
For each $w\in \V_L(v)$,  choose $\vt_w\in\Theta$ such that $\vt_w\cdot w_0=w$,
and consider
\[ m_v\coloneqq\sum_{w\in \V_L(v)}\vt_w^{-1}\cdot m_w\hs.\]
We write $\Theta_0$ for the decomposition subgroup (stabilizer)
$\Theta_{w_0}$ of $w_0$ in $\Theta$.
Then $l_v\tt\big[bas[i]\hs\big]$
is the image in $\Lambda_{\Theta_0}$ of $m_v\in\la M\ra$
under the natural homomorphism
\[ \la M\ra=\Z^\smm\to \Lambda\to\Lambda_{\Theta_0}\hs.\]
Since $ \tt x[i]\cdot bas[i]\in \la C0N\ra$, we have
\[{\tt x[i]}\cdot [m_v]=0\in \Lambda_{\Theta_0}\hs,\]
whence $[m_v]$ is a torsion element, $[m_v]\in\Lambda_{\Theta_0\hs,\Tors}$\hs.
Thus we have described the homomorphism \eqref{e:lvv}.



\section{Examples: computation of the elementary divisors of $\Sha^S(\Lambda)$}
\label{s:example}

Let $K$ be a totally real number field and let $L=K(\zeta_8)$
where $\zeta_8$ is a primitive 8th root of unity.
Since $\zeta_8^2=\sqrt{-1}\notin K$, we see that $[L:K]=4$.
Write $\Theta=\Gal(L/K)$.
We embed $\Z/8\Z$ into $(R_{L/K}\SL_8)(K)=\SL_8(L)$ by sending the generator $1+8\Z$ to the scalar matrix $\zeta_8\cdot I_8$.
We set $A=\Z/8\Z\subset R_{L/K}\SL_8$.
Following Sansuc \cite[Example 5.8]{Sansuc},  we consider  $G=(R_{L/K}\SL_8)/A$.

Write $\Lambda=\pi_1^\alg(G)$. By \cite[Corollary C.2.2]{BK}, we have
\[ \Lambda\cong\Hom(\X^*(A),\Q/\Z)\]
where $\X^*$ denotes the character group.
We obtain:
\begin{align*}
 &\X^*(A)=\Hom(\Z/8\Z,\bar K^\times)\cong\mu_8\hs,\\
&\Lambda\cong\Hom(\mu_8,\Q/\Z)=\Hom(\mu_8,\tfrac18\Z/\Z)\cong\Hom(\mu_8,\Z/8\Z).
\end{align*}
We fix an isomorphism (not $\Theta$-equivariant)
 $\Hom(\mu_8,\Z/8\Z)\isoto\Z/8\Z$
by sending a homomorphism $\chi\colon\mu_8\to\Z/8\Z$ to $\chi(\zeta_8)\in \Z/8\Z$.
We identify $\Lambda$ with $\Z/8\Z=\{\bar x=x+8\Z\ |\  x\in \Z\}$.
Then $\Aut(\Lambda)=(\Z/8\Z)^\times=\{\hs\bar 1,\bar 3,\bar 5,\bar 7\hs\}$
acting on $\Lambda=\Z/8\Z$ by  $\bar a\cdot\bar x=\ov{ax}$.
The action of $\Theta=\Gal(L/K)$ on $\Lambda$ permits us to identify  $\Theta$ with $(\Z/8\Z)^\times$.
The group $\Theta=(\Z/8\Z)^\times$ has four cyclic subgroups:
$\langle\hs\bar1\hs\rangle,\langle\hs\bar3\hs\rangle, \langle\hs\bar5\hs\rangle, \langle\hs\bar7\hs\rangle$,
and one non-cyclic subgroup $\Theta$.

By a version of \cite[Lemma 7.5.2]{BK}, for $K$ and  $G$ as above
and for each  archimedean place $v$ of $K$, the canonical map
\[H^1(K_v,G)\onto H^1_\ab(K_v,G)\into \Lambda_\Thvt\]
is in fact bijective.
Therefore, the bottom horizontal arrow in diagram \eqref{e:Cart-S} is bijective.
Now it follows from Theorem \ref{t:Sha} that the top horizontal arrow
\begin{equation}\label{e:abLambda}
\abLambda^S\colon \Sha^S(K,G)\to \Sha^S(\Lambda)
\end{equation}
in this diagram is bijective as well.
Moreover, for each  archimedean place $v$ we have $L_w/K_v\simeq\C/\R$, and therefore
the non-trivial element of $\Theta_w=\Gal(L_w/K_v)$ sends $\zeta$ to $\zeta^{-1}=\zeta^7$,
where we write $\zeta$ for $\zeta_8$.
It follows that $\Theta_v=\la\bar 7 \ra$.

We choose a finite set $\Sigma\subset \V_K$ satisfying \eqref{e:v-Sigma}.


\begin{example}\label{ex:cyclic}
We consider the case  $K=\Q\big(\sqrt{7}\hs\big)$.
By \cite[Lemma 7.5.1]{BK},
for this $K$  and the Galois extension $L=K(\zeta_8)$,
all decomposition groups are cyclic.
For each $a=1,3,5,7,$ we choose a few places   $v_a$ of $K$
with decomposition group $\langle\hs\bar a\hs\rangle$.
We take $\Sigma\smallsetminus S=\{v_3,v_5,v_7\}$; since all decomposition groups are cyclic,
condition \eqref{e:v-Sigma} is satisfied.
Using our Magma program, we compute
the finite abelian group $\Sha^S(\Lambda)$
with $S=\varnothing$, $S=\{v_a\}$ and $S=\{v_a,v_b\}$
for each $a=1,3,5,7$ and $b=1,3,5,7$.
In particular, we compute the elementary divisors of $\Sha^S(\Lambda)$,
which gives the cardinality $\#\Sha^S(K,G)=\#\Sha^S(\Lambda)$
(because the abelianization map \eqref{e:abLambda} is bijective).

%\pagebreak

\begin{longtable}{|r|l|l|l|l|l|}
\caption{$\Sha^S(\Lambda)\cong\Sha^S(K,G)$ for $K=\Q\big(\sqrt7\big)$}
\endfirsthead
\hline
\endhead
\hline
\endfoot
\endlastfoot

\hline
    &$\ \vn$    & \ $\la1\ra\hs$ &\ $\la3\ra\hs$ &\ $\la5\ra\hs$ &\ $\la7\ra\hs$
    \\
\hline
$\vn\ $   &   2    &   2,4   &   2   &   4   & 2
\\
\hline
$\ \la1\ra\ $
& 2,4 & 2,4,8 & 2,8 & 4,8 & 2,8
\\
\hline
$\ \la3\ra\ $
& 2 & 2,8 &  2,2 & 8 & 4
\\
\hline
$\ \la5\ra\ $
& 4 & 4,8 & 8 & 4,4  & 8
\\
\hline
$\ \la7\ra\ $
& 2 & 2,8  & 4  & 8 & 2,2
\\
\hline
\end{longtable}
In the cell with coordinates $\vn,\vn$ in this table,
the number $2$ is the elementary divisor
of the finite abelian group $\Sha^S(\Lambda)$
for $S=\vn$.
In the cell with coordinates $\vn, \la1\ra$,
the numbers $2,4$ are the elementary divisors for $S=\{v_1\}$.
In the cell with coordinates
$\la1\ra,\la1\ra$ ,
the numbers $2,4,8$ are the elementary divisors
of the finite abelian group $\Sha^S(\Lambda)$
for $S=\{v_1,v_1'\}$ with {\em different} places
of $K$ having trivial decomposition groups
$\la\bar 1\ra,\la\bar 1\ra$.
In the cell with coordinates $\la3\ra,\la1\ra$,
the numbers $2,8$ are the elementary divisors of the finite abelian group $\Sha^S(\Lambda)$
for  $S=\{v_3,v_1\}$ with (different) places
of $K$ having  decomposition groups $\la\bar 3\ra,\la\bar 1\ra$,
and similarly for the other cells.

In particular, for $S=\vn$ we have $\Sha^\vn(K,G)\coloneqq\Sha(K,G)\cong\Z/2\Z$,
which is well known; see \cite[(3.7)]{Sansuc} and \cite[Remark 7.5.4]{BK}.
\end{example}

\newcommand{\other}{{\rm other}}
\begin{proposition}\label{p:cubic}
Let $K=\Q(\alpha)$ where $\alpha$ is a root of the polynomial $f(x)=x^3+x^2-10x-8$.
Set $L=K(\zeta_8)$, $\Theta=\Gal(L/K)$. Then:
\begin{enumerate}
\item $K$ is totally real, with three real places.
\item $\Theta\cong(\Z/8\Z)^\times$.
\item $K$ has exactly three places $v$ with non-cyclic decomposition group,
namely, the three places over $2$ with decomposition group $\Theta$.
\end{enumerate}
\end{proposition}

We skip the proof.

\begin{example}\label{ex:Theta}
Let $K=\Q(\alpha)$ be as in Proposition \ref{p:cubic}.
Using our Magma program, we compute
the finite abelian group $\Sha^S(\Lambda)=\ker_\Sigma^S(\Lambda)$
for each $S=S_\re\cup S_\other$, where $S_\re=\{v_7^{(1)},v_7^{(2)}\}$
consists of two different {\em real} places
(with decomposition subgroup $\la\bar 7\ra$), and where
$S_\other=\varnothing$, or $S_\other=\{v_a\}$, or  $S_\other=\{v_a,v_b\}$,
for each $a=1,3,5,7,\Theta$ and $b=1,3,5,7,\Theta$.
Here by $v_a$ with $a=\Theta$
we mean a place $v_\Theta$ of $K$ with $\Theta_{v_\Theta}=\Theta$,
and similarly for $v_b$ when $b=\Theta$.
For our computation we take $\Sigma\smallsetminus S=\{\Theta_{v_\Theta}\}$;
then condition \eqref{e:v-Sigma} is clearly satisfied.
We compute the elementary divisors of $\Sha^S(\Lambda)$,
which gives the cardinality $\#\Sha^S(K,G)=\#\Sha^S(\Lambda)$
(because the abelianization map \eqref{e:abLambda} is bijective).


%\pagebreak

\begin{longtable}{|r|l|l|l|l|l|l|}
\caption{$\Sha^S(\Lambda)\cong\Sha^S(K,G)$ for $K=\Q(\alpha)$}
\endfirsthead
\hline
\endhead
\hline
\endfoot
\endlastfoot

\hline
    &$\ \vn$    & \ $\la1\ra\hs$ &\ $\la3\ra\hs$ &\ $\la5\ra\hs$ &\ $\la7\ra\hs$ & $\Theta$
    \\
\hline
$\vn\ $   &   2    &   2,8   &  2,2   &   2,4   & 2,2 & 2,2
\\
\hline
$\ \la 1 \ra\ $
& 2,8 & 2,8,8 & 2,2,8 & 2,4,8 & 2,2,8 & 2,2,8
\\
\hline
$\ \la 3 \ra\ $
& 2,2 & 2,2,8  & 2,2,2 & 2,2,4 & 2,2,2 & 2,2,2
\\
\hline
$\ \la 5 \ra\ $
& 2,4 & 2,4,8 & 2,2,4 & 2,4,4  & 2,2,4 & 2,2,4
\\
\hline
$\ \la 7 \ra\ $
& 2,2 & 2,2,8 & 2,2,2 & 2,2,4 & 2,2,2 & 2,2,2
\\
\hline
$ \Theta\ $
& 2,2  & 2,2,8 & 2,2,2 & 2,2,4 & 2,2,2 & 2,2,2
\\
\hline
\end{longtable}

The notations in Table 2 are similar to those in Table 1.
In particular, when $S_\other=\{v_1,v_5\}$,
the elementary divisors of $\Sha^S(\Lambda)$ are $2,4,8$.

\end{example}




\appendix

\section{Magma program}
\label{app:Magma}



Our program computes the relative Tate--Shafarevich group
$\ker_{\Sigma}^S(\Theta, \Lambda)$ with $\Lambda=\la M\ra/\la N\ra $.
Here $\Theta$ is a finite group of order $\tt gg$;
it is denoted by $G$ in the program.
Moreover, $\tt\la M\ra= \Z^\smm$ is a free
finitely generated abelian group with $\smm$ generators,
and $\tt\la N\ra$ is a subgroup of $\Z^\smm$ with $\tt nn$ generators
that are the rows of an $\tt nn\times \smm$ matrix $\tt N$.

Furthermore, $\Sigma= S\cup \Sigma\smallsetminus S$
is a  finite sequence $(\Thv)_{v\in\Sigma}$ of {\em  subgroups of $\Theta$}.
The result depends only on the conjugacy classes of the subgroups.
Actually, the result depends only on the multiplicity of each conjugacy class.
Moreover, the result does not depend
on the positive multiplicities for $\Sigma\smallsetminus S$,
and so we assume that all these multiplicities are $1$ (or $0$).
Returning to our field extension $L/K$,
we assume that \eqref{e:v-Sigma} holds.


In addition to $\ker_{\Sigma}^S(\Theta, \Lambda)$, our program computes also
the homomorphisms
\[l_v\colon \ker_{\Sigma}^S(\Theta, \Lambda) \to \Lambda_{\Theta_v,\Tors}\]
for $v\in S_\re$; see \eqref{e:lvv}.
Here $S_\re$ is the subsequence of $S$ corresponding to the real places of $K$;
then $\#\Theta_v=2$ for all $v\in S_\re$.


For explicit calculations we chose an example with $\smm=1$.
We take $\tt \la M\ra=\Z$, $\tt nn=1$, $\tt N=[[8]]$;
then $\tt\la N\ra= 8\Z$ and  $\Lambda=\Z/8\Z$.
We take $\Theta=\Aut \Lambda=(\Z/8\Z)^\times=\{\bar 1,\bar3,\bar5,\bar 7\}$
where for $a\in\Z$ we denote  $\bar a=a+8\Z$, the image of $a$ in $\Z/8\Z$.





\bigskip

\centerline{\sc The program}
\medskip

\begin{lstlisting}[language=C++]

/* This program computes Sha^S(G,M/N)
 where G is a finite group of order gg
acting on a finitely generated abelian group M/N,
where M is the free abelian group  Z^mm
with mm generators,
and N is a subgroup of M with nn generators.
*/


ComplementIntMat:= function( full, sub )
/* This is a Magma version, written by Willem A. de Graaf,
of the GAP function with the same name.
*/
     full:= Matrix( Integers(), full );
     H:= HermiteForm( full );
     full:= [ ];
     for i in [1..NumberOfRows(H)] do
         if not IsZero( H[i] ) then
            Append( ~full, Eltseq( H[i] ) );
         end if;
     end for;
     full:= Matrix( Integers(), full );

     n:= NumberOfColumns(full);
     m:= NumberOfRows( full );
     Vn:=KModule( Rationals(), n );
     W:= KModuleWithBasis( [ Vn!Eltseq(full[i]) : i in [1..m] ] );
     cfs:= [ Coordinates( W, Vn!Eltseq(sub[i]) ) : i in [1..NumberOfRows(sub)] ];
     A:= Matrix( Integers(), cfs );

     S,P,Q:= SmithForm( A );
     Qi:= Q^-1;
     cfs:= [ Eltseq(Qi[j]) : j in [1..m] ];
     vecs:= [ &+[ c[i]*(Vn!Eltseq(full[i])) :
        i in [1..m] ] : c in cfs ];
     x:= [ ];
     minRC:=Minimum(NumberOfRows(S),NumberOfColumns(S));
     for i in [1..minRC] do
         if S[i][i] ne 0 then
            Append( ~x, S[i][i] );
         end if;
     end for;

     comp:= [ ];
     for i in [#x+1..#vecs] do
        Append( ~comp, Eltseq(vecs[i]) );
     end for;

     return Matrix( Integers(), comp),
	 Matrix( Integers(),
     // [ Eltseq(x[i]*vecs[i]) : i in [1..#x] ] ),
	 	[ Eltseq(vecs[i]) : i in [1..#x] ] ),
	    Vector( Integers(), x );

end function;



IsGroup := function(gg, mu)
/*This function checks that the pair (G,mu),
where G={1,..,gg} and mu: G x G -> G,
is a group, and if yes,
computes the inversion map iota.
*/
	boo:=true;	
	
	for g:=1 to gg do
		boo := boo and (mu[g,1] eq g) and (mu[1,g] eq g);
//      the unit condition
	end for;	
				
	for g1:=1 to gg do
	for g2:=1 to gg do
		for g3:=1 to gg do
			p12:=mu[g1,g2]; p23:=mu[g2,g3];
			boo:=boo and (mu[p12,g3] eq mu[g1,p23]);
//			 the associativity condition
		end for;
	end for;
	end for;
		
	iota:=[];

	for g:=1 to gg do
		boog:=false;
		for h:=1 to gg do
			boo1:=(mu[g][h] eq 1);
			if boo1 then
		    	iota[g]:=h;
				boog:=true;
				break;
			end if;
		end for;
		boo:=boo and boog;
//		 existence of inverse
	end for;
		
	return boo,iota;
	
end function;



IsAction:= function(gg,mu,mm,rho,nn,N)
/* This function, for a matrix N with mm columns
and a family \rho(g) of mm x mm matices,
checks whether rho induces a left action
of (G,mu) on Z^mm / <N> .
*/	
	boo:=true;	
	Zv:=RModule(Integers(), mm);
	vv:= Zero(Zv);			

	for g:=1 to gg do
	for i:=1 to nn do
		for j:=1 to mm do
			vv[j]:=N[i][j];
		end for;
		vv:=vv * Transpose(rho[g]);	
		ic:=IsConsistent(N, vv);
//		rho preserves N		
		boo:=boo and ic;
	end for;
	end for;
	
	Idm := ScalarMatrix(Integers(), mm, 1 );
	Dif:= Transpose (rho[1]-Idm);
	for i:=1 to mm do
			for j:=1 to mm do
				vv[j]:=Dif[i][j];
			end for;	
			ic:=IsConsistent(N, vv);
// 			rho[1] is 1  modulo N
			boo:=boo and ic;
		end for;
	
	for g:=1 to gg do
	for h:= 1 to gg do
		Dif:=Transpose( rho[mu[g][h]] - rho[g]*rho[h] );
		for i:=1 to mm do
			for j:=1 to mm do
				vv[j]:=Dif[i][j];
			end for;	
			ic:=IsConsistent(N, vv);
// 			rho is a well-defined left action modulo N
			boo:=boo and ic;
		end for;
	end for;
	end for;
		
	return boo;
	
end function;
			


Meet := function(A,B) //
/* This function computes a matrix R
such that <R> = <A> \cap <B>.
*/	
	AB:=VerticalJoin(A,B);
	NS:=NullspaceMatrix(AB);
	rA:=NumberOfRows(A);
	rNS:=NumberOfRows(NS);
	R:=ExtractBlockRange(NS, 1, 1, rNS,rA) * A;
	
	return R;
	
end function;
	
	
	
Mult:= function(g,X,mu)
// Left action of an element g of G
// on a subset X of G.	
	return { mu[g][x] : x in X };
end function;



CheckSubGroup:= function(G,H,mu,iota)
	
	boo:=true;
	for h1 in H do
		boo1:=iota[h1] in H;
	if not boo1 then
		boo:=false; break;
	else
		for h2 in H do
			boo2:= mu[h1][h2] in H;
			if not boo2 then
				boo:=false; break;
			end if;
		end for;
	end if;
	end for;
	
	return boo;
	
end function;



Cosets:=function(gg,H,mu)
/* This function computes the left cosets
of G={1..gg} by a subgroup H.
*/
	cosets:=[]; ic:=0;
	booG:=[true : g in [1..gg] ];
	
for g:=1 to gg do
	if booG[g] then
		ic:=ic+1;
		cosets[ic]:= Mult(g,H,mu);
		for x in cosets[ic] do
			booG[x]:=false;
		end for;
	end if;
end for;

return cosets;

end function;




Num:= function(X, cosets, yy)
// The number (sic!) of a given left coset X of G.

	for i:= 1 to yy do
		if X eq cosets[i] then
			numX:=i; break;
		end if;		
	end for;
	
	return numX;
	
end function;
	
	
	


NotOne:= function(bas,x)

/* This function computes the sequence nOne of the elements
of the integer vector x different from 1,
and the matrix of the corresponding rows of the matrix bas
*/
	lbas:=NumberOfColumns(bas);
	d:=#Eltseq(x);
	gen:=[];
	s:=[];
	l:=0;
	
	for i:=1 to d do
		if x[i] ne 1 then
			Append(~s,x[i]);
			l := l+1;
			Append(~gen,bas[i]);
		end if;
	end for;
	
	gen := Matrix(Integers(), l, lbas,  gen);
	
	return  s, gen;
end function;




// ######################################

print "\n\nINPUT";	

//INPUT: G (that is, mu)
// two integers gg, mm.
	
gg:=4;
	
//  A list of lists mu[1..gg][1..gg]
// (the multiplication law)
	 mu:=[[1,2,3,4],[2,1,4,3],[3,4,1,2],[4,3,2,1]];		

// INPUT of Lambda, that is, mm, N, and rho
	 	
mm:=1;	

		
// #######################

rho:= [ ZeroMatrix(Integers(),mm,mm) :
    u in [1..gg] ];

rho[1]:=[[1]];
rho[2]:=[[3]];
rho[3]:=[[5]];
rho[4]:=[[7]];

//	INPUT: N
	
nn:=1;

N:=ZeroMatrix(Integers(),nn, mm);
N[1][1]:=8;


print "\ngg=", gg,
    "\nthe multiplication matrix: mu\n", mu;	
	
print "\nmm=", mm, "<M> = <I_mm>= Z^mm";

print "\nthe action of G: rho \n", rho;	

print "\nnn=", nn,
    "\nthe relations in <M> = Z^mm:  N", N;

print "\nOur G-module is Z^mm / <N>";	
	
boo, iota := IsGroup(gg,mu);	
print "\nIs G a group?  ", boo;
print "\nthe inversion map: iota ", iota;
		
print "\nIs rho an action of G on Z^mm / <N> ?  ",
    IsAction(gg,mu,mm,rho,nn, N);




// ##################################


print "\n\n\nCOMPUTING Sha^S(K,G)";




// G = {1... gg}

// The subgroups of G: HH[i]
HH:=[];

HH[1]:={1};
HH[2]:={1,2};
HH[3]:={1,3};
HH[4]:={1,4};
HH[5]:={1,2,3,4};

mHH := [];
mHH[1]:=[1];
mHH[2]:=[1,3];
mHH[3]:=[1,5];
mHH[4]:=[1,7];
mHH[5]:=[1,3,5,7];

boo := true;
for HHH in HH do
	print "HHH = ", HHH;
	boo := boo and  CheckSubGroup( gg, HHH, mu, iota );
end for;
print "All subsets HH[i] are subgroups: ", boo;





ShaS:= function(gg, mu, iota, mm, N, rho, HH,
    k_1re, k_1other, k_2 )

/* This is the main function of the program.
It computs Sha^S= ker_\Sigma^S. Here
k_1re is the sequence of real places in S,
k_1other is the sequence of other places in S, and
k_2 is the set of places in  Sigma - S.
*/



nn:= NumberOfRows(N);


k_1 := k_1re cat k_1other;

kk_1:=#k_1;
kk_2:=#k_2;


k:=k_1 cat k_2; // Sigma
kk:=#k;  // The number of places in Sigma


xx:=[ 0 : u in [1..kk]];

for i:=1 to kk do
	xx[i]:=gg div #HH[k[i]];
    // The number of places of L
    // over the place v[i] of K.
end for;


// #########################

//Computing the matrices A[g][i] of action of g
// on the left cosets modulo HH[k[i]]

AA:=[ [* ZeroMatrix(Integers(), xx[i], xx[i]) :
    i in [1..kk] *] : g in [1..gg] ];

for i:=1 to kk do
for g:=1 to gg do
	AA[g][i]:=ZeroMatrix(Integers(), xx[i], xx[i]);
	cosets:=Cosets(gg,HH[k[i]],mu);
	for x:=1 to xx[i] do
//		print g, i, x, xx[i];
		y:=Num(Mult(g,cosets[x],mu), cosets, xx[i]);
		AA[g][i][x][y]:=1;
	end for;
end for;
end for;


/* Here AA[g][i] is the permutation matrix
of the action of g on the places of L
lying over the place v[i] of K.
*/
	
// ###############################

// Computing the matrices Z0 and C0 for i:=1 to kk
//for zero-sum vectors

Z0:=[];
zz:=0;

C0:=[];
cc:=0;

lx:=0;
dd:=0;

for i:=1 to kk do
for x:= 1 to xx[i] do
	lx:=lx+1;	
for m:=1 to mm do
	dd:=dd+1;	
		
	B:=[* ZeroMatrix(Integers(), xx[i], mm) : i in [1..kk] *];
	Bp:=[* ZeroMatrix(Integers(), xx[i], mm) : i in [1..kk] *];

	B[i][x][m]:=1;

	if (i eq 1) and (x eq 1) then
		B[1][1][m]:=0;
	else
		B[1][1][m]:=-1;
	end if;
	
	z0:=&cat[ElementToSequence(B[j]) : j in [1..kk] ];
	Append(~Z0,z0); // The matrix of generators of the group
    // of zero-sum vectors in k_1 and k_2
	zz:=zz+1;
	
	for g:=1 to gg do
		Bp:=B;		
		Bp[i]:=B[i]-Transpose(AA[g][i])*B[i]*Transpose(rho[g]);
		Bp[1]:=B[1]-Transpose(AA[g][1])*B[1]*Transpose(rho[g]);
		b:=&cat[ElementToSequence(Bp[j]) : j in [1..kk] ];
		Append(~C0,b);	// The matrix of x - g.x
        // for zero-sum vectors x	in k_1 and k_2
		cc:=cc+1;
	end for;
	
end for;
end for;
end for;
	
	
Z0Mat:=Matrix(Integers(), zz, dd, Z0 );
C0Mat:=Matrix(Integers(), cc, dd, C0 );
// Two matrices with dd = mm * kk columns


for l:=1 to lx do
	if l eq 1 then
		NN:=N;
	else
		NN:=DiagonalJoin(NN,N);
	end if;
end for;
// NN is a matrix with dd columns: <NN> = <N>[Sigma]

Z0N:=VerticalJoin(Z0Mat,NN); // <Z0N> = Z0Mat> + <NN>
C0N:=VerticalJoin(C0Mat,NN); // <C0N> = C0Mat> + <NN>



// ################################################

// Computing the matrix CC for S for i:=1 to kk_1

CC:=[];

lx:=0; dd:=0; cc:=0;

for i:=1 to kk_1 do
for x:= 1 to xx[i] do
	lx:=lx+1;	
	for m:=1 to mm do
		dd:=dd+1;	
		
		B:=[* ZeroMatrix(Integers(), xx[i], mm) :
            i in [1..kk] *];
		Bp:=[* ZeroMatrix(Integers(), xx[i], mm) :
            i in [1..kk] *];

		Bp[i][x][m]:=1;
		b:=&cat[ElementToSequence(Bp[j]) : j in [1..kk] ];
		Append(~CC,b);
            //The matrix of ALL vectors in S (k_1).	
		cc:=cc+1;
	end for;
end for;
end for;


// ############################################

// Computing CC for Sigma - S
// for i:= kk_1+1 to kk_1+kk_2

for i:=kk_1+1 to kk_1+kk_2 do
for x:= 1 to xx[i] do
	lx:=lx+1;	
for m:=1 to mm do
	dd:=dd+1;	
		
	B:=[* ZeroMatrix(Integers(), xx[i], mm) :
        i in [1..kk] *];
	Bp:=[* ZeroMatrix(Integers(), xx[i], mm) :
        i in [1..kk] *];

	B[i][x][m]:=1;
	
	for g:=1 to gg do
		Bp:=B;
		Bp[i]:=B[i]-Transpose(AA[g][i])*B[i]*Transpose(rho[g]);
		b:=&cat[ElementToSequence(Bp[j]) : j in [1..kk] ];
		Append(~CC,b);	// The matrix of x - g.x
             // for all vectors x in k_2.
		cc:=cc+1;
	end for;
end for;
end for;
end for;


CCMat:=Matrix(Integers(), cc, dd, CC );	
CCN:=VerticalJoin(CCMat,NN);
/*
CCN is the set of all cocycles (=cochains)  in M[Sigma] (modulo <N>)
( where Sigma = k_1 \cup k_2 )
that become coboundaties over Sigma - S = k_2
*/


// ####################################

m_Z0N_CCN := Meet(Z0N,CCN);
// <m_Z0N_CCN> = <Z0N> \cap <CCN>
comp, bas, x:=ComplementIntMat( m_Z0N_CCN, C0N );
/* x =
=  ( ( <Z_0N> \cap <CCN> ) / <C0N> )_Tors
= ker_\Sigma^S
*/



if k_2 eq [2,3,4] then
sk_2:= "CYCLIC: [ [1,3], [1,5], [1,7] ]";
else sk_2:="THETA: [ [ 1, 3, 5, 7 ] ]";
end if;

// Printing the list of the decomposition
// subgroups in k_1re,  k_1other, and k_2.

Pk_1re:= "[ ";
for i in k_1re do
	Pk_1re:= Pk_1re cat Sprint(mHH[i]) cat "  ";
end for;
Pk_1re:=Pk_1re cat "]";

Pk_1other:= "[ ";
for i in k_1other do
	Pk_1other:= Pk_1other cat Sprint(mHH[i]) cat "  ";
end for;
Pk_1other:=Pk_1other cat "]";


print "\nk_1re:", Pk_1re,  "\nk_1other:", Pk_1other,
"\nS = k_1 := k_1re cat k_1other";
// [ mHH[i] : i in k_1 ];
print "Sig - S = k_2:", sk_2, "\n";
// [ mHH[i] : i in k_2 ];

nOne, gen:= NotOne(bas,x);

return nOne, gen, xx;

end function;

//###########################################



ellV := function(gg, mu, iota, mm, N, rho, HH, k_1re, l, gen , xx)

/* This function computes the homomorphisms l[v]
for the real places in v in S_\infty = k_1re.
*/
	if l eq 0 then return []; end if;
	
	mm_ellV := [];
	kk_1re := #k_1re;
	lv:=0;
	lw:=0;
	
	for v := 1 to kk_1re do
		lv:=lv+1;
		matV := [];
		for i := 1 to l do
			cosets := Cosets(gg, HH[k_1re[v]], mu );
			lw:=0;
			for vv:= 1 to v-1 do
				lw := lw+xx[vv];
			end for;
			
			for y := 1 to xx[v] do
				lw := lw+1;
				vecy:=ZeroMatrix(Integers(), 1, mm);
				for g := 1 to gg do
					yNum :=
                            Num( Mult( g, HH[k_1re[v]], mu ),
                            cosets, xx[v]);
					if yNum eq y then
						gy := g; break;
					end if;
				end for ; // g	
				
			vec:=Matrix(Integers(),1, mm,
			[ gen[i][ (lw-1)*mm+j ] : j in [1..mm] ]);
			vecy := vecy+ vec*Transpose(rho[iota[gy]]);
			end for;  // y
		
			Append( ~matV, Eltseq(vecy) );
		end for; // i
		Append( ~mm_ellV, matV);
	end for; // v
	R_ellV:=mm_ellV;
	
	if mm eq 1 then
		R1 := ZeroMatrix( Integers(), l, kk_1re );
		
		for v := 1 to kk_1re do
			for i := 1 to l do
				R1[i][v] :=mm_ellV[v][i][1];
			end for;
		end for;
		R_ellV:=R1;
	end if;
		
	return R_ellV;
	
end function;
				
	


Running := procedure(k_1re, k_1other, k_2)

// This procedure runs the main functions ShaS
// and ellV and prints the results.

	nOne, gen, xx := ShaS(gg, mu, iota, mm, N, rho, HH,
        k_1re, k_1other, k_2 );
	print "nOne=", nOne;	
	l := #nOne;
	print "gen = ",  gen;
	
    if k_1re ne [] then
    ellvE := [];
    ellvE :=  ellV( gg, mu, iota, mm, N, rho, HH, k_1re, l, gen , xx );
    for i:= 1 to NumberOfRows(ellvE) do
    	for j:= 1 to NumberOfColumns(ellvE) do
    		ellvE[i][j] := ellvE[i][j] mod 2 ;
    	end for;
    end for;
	print "ellV =", ellvE;
    end if;
	print "*********************";
end procedure;




//##########################
// THE MAIN CYCLES

// (1) : Below we compute Sha^S(Lambda) for K = Q(\sqrt{7}) and $L=K(\zeta_8)
// as in Example 5.1.

print "\nCase of  K = Q(sqrt{7}) \n\n";


k_1re := [];  // The set of *real* places in S
kk_1re := 0;


k_2:= [2,3,4];
sk_2:= "CYCLIC";
	
		

	print "########################################\n";
	print "## k_1 is EMPTY,  k_2 is ", sk_2, "  ##";
	
	k_1other:= [];
	Running(k_1re, k_1other, k_2);
	

	print "#########################################\n";
	print "## k_1 contains ONE place,  k_2 is ", sk_2, "  ##";
	
	for ii := 1 to 4 do
		
		k_1other := [ii];
		Running(k_1re, k_1other, k_2);
		
	end for;
	

	print "#############################################\n";
	print "## k_1 contains TWO places,  k_2 is ", sk_2, "  ##";

	for ii := 1 to 4 do
	for jj := 1 to 4 do
	
		k_1other := [ii,jj];
		Running(k_1re, k_1other, k_2);
		
	end for;
	end for;


/* (2) : Below we compute Sha^S(Lambda) for K = Q(alpha) and $L=K(zeta_8)
 as in Example 5.3. Here alpha  is a root of the polynomial
f(x) = x^3 + x^2 - 10*x - 8.  */

print "\nCase of  K = Q(alpha)  where  alpha  is a root \nof the polynomial",
"f(x) = x^3 + x^2 - 10*x - 8.\n\n";

k_1re := [ 4, 4 ];  //  The sequence of real places in S
kk_1re := #k_1re;

print "k_1re contains two places: ", k_1re;

booRe:= true;
for v in k_1re do
//	print "v = ", v, "   HH[v] = ", HH[v];
  	if #HH[v] ne 2  then
  		print "HH[v] is not real for v = ", v;
  		booRe := false;
  		break;
  	end if;
end for;
print "All subgroups  HH[v] for v in k_1re are of order 2: ", booRe;

k_2 := [ 5 ];
sk_2:="THETA";
		

	print "########################################\n";
	print "## k_1other is EMPTY,  k_2 is ", sk_2, "  ##";
	
	k_1other:= [];
	Running(k_1re, k_1other, k_2);
	

	print "#########################################\n";
	print "## k_1other contains ONE place,  k_2 is  ", sk_2, "  ##";
	
	for ii := 1 to 5 do
		
		k_1other := [ii];
		Running(k_1re, k_1other, k_2);
		
	end for;
	

	print "#############################################\n";
	print "## k_1other contains TWO places,  k_2 is ", sk_2, "  ##";

	for ii := 1 to 5 do
	for jj := 1 to 5 do
	
		k_1other := [ii,jj];
		Running(k_1re, k_1other, k_2);
		
	end for;
	end for;


\end{lstlisting}



\newcommand{\sk}{\smallskip}

\begin{thebibliography}{BdG24}

\bibitem[Brl63]
{Borel}
Armand Borel,
{\em Some finiteness properties of adele groups over number fields},
Inst. Hautes Études Sci. Publ. Math. No. {\bf 16} (1963), 5--30.
\sk

\bibitem[Brv98]
{Borovoi-Memoir}
Mikhail Borovoi,
{\em Abelian Galois cohomology of reductive groups,}
Mem. Amer. Math. Soc. {\bf 132} (1998), no. 626.
\sk

\bibitem[BdG24]
{BdG}
Mikhail Borovoi and Willem A. de Graaf,
{\em Computing Galois cohomology of a real linear algebraic group},
J. Lond. Math. Soc., II. Ser., {\bf 109}, No. 5, Article ID e12906, 53 p. (2024).
\sk

\bibitem[BK26]
{BK}
Mikhail Borovoi and Tasho Kaletha,
with an appendix by Vladimir Hinich,
{\em Galois cohomology of reductive groups over global fields,}
Mem. Amer. Math. Soc. {\bf 319} (2026), no. 1621, vii+110 pp.
\sk

\bibitem[BT24]
{BT}
Mikhail Borovoi and Dmitry A. Timashev,
{\em Galois cohomology and component group of a real reductive group},
Israel J. Math. {\bf 259} (2024), no. 2, 937--1000.
\sk

\bibitem[CT08]
{CT}
Jean-Louis Colliot-Th\'el\`ene,
{\em R\'esolutions flasques des groupes lin\'eaires connexes},
J. Reine Angew. Math. {\bf 618} (2008), 77--133.

\bibitem[Ko84]
{Kottwitz-84}
Robert E. Kottwitz,
{\em Stable trace formula: cuspidal tempered terms},
Duke Math. J. {\bf 51} (1984), no. 3, 611--650.
\sk

\bibitem[Rot09]
{Rotman}
 Joseph J. Rotman,
{\em An introduction to homological algebra,
second edition},
Universitext,
Springer, New York, 2009.
\sk

\bibitem[San81]
{Sansuc}
Jean-Jacques Sansuc, \emph{Groupe de {B}rauer et arithm\'etique des groupes
alg\'ebriques lin\'eaires sur un corps de nombres}, J. Reine Angew. Math.
{\bf 327} (1981), 12--80.

\end{thebibliography}

\end{document}


